\subsection*{电流图像中容易混淆的三个时间尺度}

导体与绝缘体的区别在于，导体内部存在能够在外场作用下产生宏观迁移的载流子。金属、电解质、半导体、
电离气体和超导体虽然都能导电，内部机制却截然不同。金属中有近似离域的电子；电解质中正负离子向相反
方向漂移却可能产生同向的传统电流；半导体的载流子浓度会受温度、掺杂和光照显著影响；气体通常要先被
电离才有足够载流子；超导态更不能用普通散射模型解释。把这些材料都写成一个电阻符号，是电路层次的抽象，
并不意味着它们在微观层次遵循同一种过程。

经典电子模型的局限需要在开始时说明。金属电子不能用经典能量均分定理确定速度，它们服从量子统计；周期
晶格中的运动也必须用能带描述。这里使用漂移速度、平均自由时间和有效碰撞，是为了提取与宏观电流直接相关
的统计量。模型只要能够解释 $j\propto E$、Joule 热、暂态响应和温度趋势，就已经完成本讲的任务；它不能
替代近代固体物理理论。

银导线中电子的漂移速度只有 $10^{-4}\,\si{m.s^{-1}}$ 的数量级，不能据此认为合上开关后，要等电子从电源
缓慢漂移到灯泡，灯才会亮。这里存在三个完全不同的时间尺度。第一是电磁场和表面电荷重新分布的传播时间，
信号沿电路以接近电磁波在周围介质中的传播速度建立；第二是电子平均漂移速度发生改变所需的弛豫时间，典型
金属中可短到 $10^{-14}\,\si{s}$；第三才是某个电子在很长导线中产生可观净位移所需的漂移时间。电路快速
响应由前两个尺度决定，不是靠同一批电子从电源跑到负载。

电流定义中的“通过截面”也必须带方向。若正电荷沿截面法向通过，记为正电流；负电荷反向通过同样产生正
电流。对于电子导电，电子速度与传统电流方向相反，但 $q<0$，所以 $\vect j=nq\vect v_d$ 仍自动给出正确
方向。多载流子系统中不能简单比较粒子速度，还要把每类载流子的电荷符号与浓度一起计入。

稳恒电流并不要求每个电子做匀速直线运动，也不要求任一点没有电荷。它只要求宏观电流分布不随时间改变，
局部电荷密度也不再持续积累。导线表面可以存在极小而稳定的电荷梯度，正是这些表面电荷建立了导线内部所需
的电场。若在某节点流入电流大于流出电流，节点电荷会增加并改变周围电势，直到新的电场把各支路电流调整到
满足 KCL 的状态。

\subsection*{从端口观点理解元件}

把元件画成方框并不是说内部结构不重要，而是根据研究目的选择描述层次。分析收音机电路时，可能只需要晶体
管的端口特性；设计晶体管时，则必须研究载流子的注入、复合和能带。端口模型的价值在于模块化：只要每个元件
的外部关系和连接拓扑已知，就可以不重复求解内部所有自由度。

规定端口电压、电流方向后，伏安图必须明确坐标轴。若横轴为 $U$、纵轴为 $I$，过原点直线的斜率是
$I/U=1/R$；若交换坐标轴，斜率才是 $R$。同一工作点处还可以区分静态电阻 $U/I$ 与微分电阻
$\dd U/\dd I$。非线性元件的这两个量通常不同，负微分电阻表示局部斜率为负，并不等于元件在该点的
$U/I$ 必然为负。

理想电阻的伏安曲线经过原点，正、负电压下具有反向对称性。实际金属电阻通电后会发热，温度改变又使电阻率
改变，若测量过程不能保持温度恒定，伏安曲线也会弯曲。因此，Ohm 定律中的“外部状态不变”不是可有可无的
附加语，而是把 $R$ 当作常量的必要条件。

元件功率的正负由参考方向共同决定。采用被动符号约定时，电流从电压正端流入，$P=UI>0$ 表示电能进入元件；
电源正常供能时，实际电流常从其正极流出，相对于被动约定 $I<0$，因此 $P<0$。只要参考方向定义清楚，
不需要为电阻和电源分别发明两套功率公式。

静态伏安曲线还要求测量过程足够慢，使元件内部状态已经稳定。电容器在直流稳态下电流为零，却在充电时有
电流；电感在直流稳态下近似导线，却在电流变化时产生感应电压。若只给出某一瞬间的 $U,I$，不能判断这些
储能元件的后续行为。一般电路理论之所以要把电阻、电容、电感和受控源分开，就是因为它们需要不同的状态
关系。当前章节先集中研究不储存可恢复电磁能的电阻型元件。

\subsection*{从宏观电阻到微观电导率}

$R=\rho l/A$ 把材料性质与几何性质分开。将同一种均匀电阻丝长度加倍，相当于把两个相同电阻串联，电阻
加倍；将两根相同电阻丝并排连接，相当于截面积加倍，电阻减半。这一几何推理与连续介质计算完全一致。
电阻率 $\rho$ 才是给定温度和材料状态下更接近材料本身的物理量。

局部 Ohm 定律可以从圆柱导体反推。均匀导体内 $E=U/l$、$j=I/A$，代入
$U=I\rho l/A$，消去 $l,A$ 得 $E=\rho j$。各向同性材料中 $\vect j$ 与 $\vect E$ 同向，写成
$\vect j=\sigma\vect E$。若材料各向异性，电导率应当是张量，电流密度与电场不一定平行；若材料不均匀，
$\sigma$ 还会随位置变化。

体积功率密度 $p_V=\vect j\cdot\vect E$ 可由单个电荷的功率相加得到。单位体积内第 $s$ 类载流子的受力
功率为 $n_sq_s\vect v_{d,s}\cdot\vect E$，对全部载流子求和正是 $\vect j\cdot\vect E$。在线性被动介质中
$p_V=\sigma E^2\geq0$，这些能量最终通过散射传给晶格，表现为 Joule 热。

\subsection*{两种经典导电模型的联系}

黏滞阻力模型写作 $m\dot{\vect v}_d=q\vect E-\alpha\vect v_d$。与 Drude 形式比较可知
$\alpha=m/\tau$，两种写法只是参数选择不同。恒定电场下方程的完整解为
\begin{equation}
 \vect v_d(t)=\vect v_{\infty}+
 \left[\vect v_d(0)-\vect v_{\infty}\right]\e^{-t/\tau},
 \qquad
 \vect v_{\infty}=\frac{q\tau}{m}\vect E.
\end{equation}
它同时包含电流建立和撤场后的衰减。若外场变化周期远大于 $\tau$，电子几乎总能跟上瞬时稳态，静态 Ohm
定律是好近似；若频率高到与 $1/\tau$ 同量级，就必须保留惯性项，电导率变成与频率有关的复数量。

用平均自由飞行图像估计漂移速度时，电子在一次碰撞后若完全失去定向动量，随后在电场中从零加速。一次自由
飞行中的平均定向速度严格含有约 $1/2$ 的因子；但碰撞时间有统计分布，简单数量级估计常把这个系数吸收到
$\tau$ 的定义中。竞赛推导应先说明所用模型，再保持同一套定义，避免一处用平均时间、一处用固定碰撞时间
而造成因子不一致。

经典模型把碰撞画成电子与固定原子小球相撞，这只是示意。周期晶格本身并不会像随机障碍物那样阻止理想 Bloch
电子运动，真正产生电阻的是晶格振动、杂质、缺陷和边界等破坏平移对称性的散射机制。银的平均自由程远大于
原子间距，正是朴素“小球连续碰撞”图像失效的证据之一。

\subsection*{温度、场强和载流子浓度}

电导率 $\sigma=nq^2\tau/m$ 中的每个量都可能依赖外部条件。普通金属的导电电子浓度主要由价电子数和原子
密度决定，在不发生相变的温区变化较小；升温主要增强声子散射，使 $\tau$ 缩短，电阻率上升。常温附近常用
近似
\begin{equation}
 R(T)=R_0\left[1+\alpha_R(T-T_0)\right]
\end{equation}
描述，但温区很宽时必须使用非线性关系。

本征半导体中的载流子需要由热激发跨越能隙产生，$n(T)$ 对温度极为敏感。升温虽然也会缩短散射时间，却可能
使载流子数量增加得更快，因而总电导率上升。掺杂半导体还会经历杂质电离、饱和和本征激发等不同温区，不能
用一句“半导体电阻随温度下降”覆盖全部情况。

强电场会使电子在一次自由飞行中获得不可忽略的能量，改变碰撞截面、迁移率甚至载流子浓度。气体中电子可
通过碰撞电离产生更多电子和离子，形成雪崩；半导体中也可能出现速度饱和、碰撞电离和击穿。此时 $n$ 与
$\tau$ 都随场改变，线性 Ohm 定律自然失效。

气体放电曲线上不同分支对应不同放电机制和空间结构。同一电压可能对应不止一个稳定或暂稳电流，元件状态还
取决于过去的放电历史。分析这类元件时，不能仅用一个静态电阻值替代完整伏安曲线，更不能把曲线某段的斜率
直接当成始终有效的电阻。

Ohm 定律是唯象定律而不是自然界无条件成立的基本定律。它的地位类似小形变下的 Hooke 定律：在适用范围内
非常准确，超出范围后必须改用更一般的材料关系。工程师频繁使用 Ohm 定律，恰恰是因为许多金属元件在正常
温度与电流范围内足够接近线性；这并不能推出所有导电现象都应当线性。判断题目能否使用 $U=IR$ 时，应先
检查材料、温度、场强和工作状态是否允许把 $R$ 当作常量。

\subsection*{网络方程为什么会有冗余}

对含 $N$ 个节点的连通网络逐节点列 KCL，会得到 $N$ 个方程；但每一支路电流在两个端点方程中分别以相反
符号出现，把全部节点方程相加后左边恒为零。因此独立 KCL 只有 $N-1$ 个。少掉的一个方程并不是信息丢失，
而是整个网络总电荷守恒已经包含在其余节点方程中。

KVL 的冗余可用回路叠加理解。如果回路 $L_3$ 的有向边恰好等于回路 $L_1$ 与 $L_2$ 的有向边之和，那么
$L_3$ 的电压方程就是前两个方程相加。对于有 $B$ 条支路、$N$ 个节点的连通网络，独立回路数为
$B-N+1$。选择一棵包含全部节点而没有回路的生成树，每增加一条树外支路就产生一个独立基本回路。

KCL 给出 $N-1$ 个独立约束，KVL 给出 $B-N+1$ 个独立约束，两者合计 $B$ 个，再配合 $B$ 条支路各自的
伏安关系，恰好构成求解全部支路电流和电压的完备方程组。这一计数解释了为什么既不能少列必要方程，也没有
必要把图中所有能看到的回路都列一遍。

KCL 与 KVL 的适用也包含集总近似。电路尺寸远小于信号变化对应的电磁波波长时，可以把导线近似为瞬时等势
连接，把电场和磁场效应集中到元件参数中。若电路很大或频率很高，电压、电流会沿传输线传播，同一根导线的
不同位置不再处于同一瞬间状态，简单节点和支路模型需要升级为分布参数模型。竞赛中的直流网络通常默认满足
集总条件，但理解这一前提能防止把 KVL 错用于穿过显著时变磁通的回路。

\subsection*{节点电压法的列式模板}

节点法处理纯电阻支路时，可固定参考节点电势为零。若未知节点 $i$ 与相邻节点 $j$ 通过电阻 $R_{ij}$ 相连，
从 $i$ 流向 $j$ 的电流写成
\begin{equation}
 I_{ij}=\frac{\varphi_i-\varphi_j}{R_{ij}}.
\end{equation}
对节点 $i$ 列 KCL：
\begin{equation}
 \sum_j\frac{\varphi_i-\varphi_j}{R_{ij}}=I_i^{\rm inj},
\end{equation}
其中右侧表示外部电流源注入该节点的净电流。所有未知节点都得到一个线性方程，系数矩阵的对角元是与该节点
相连的电导之和，非对角元是相邻节点间电导的负值。

若理想电压源连接未知节点与参考地，该节点电势直接由电源确定；若理想电压源跨接两个未知节点，可把两节点
组成超节点，对超节点整体列 KCL，再增加电压约束。节点法并没有废除 KVL，而是用节点电势差自动保证所有
不含时变磁通的回路满足 KVL。
